\footnotesize%
\pstart%
\noindent%
Die Aufzeichnungen N.\;\ref{57040} und \ref{57041} weisen den gleichen Titel -- \textit{Motus perpetuus per thermometrum, barometrum, hygrometrum} -- auf und bilden einen einheitlichen, bruchlosen Text, in welchem unterschiedliche mechanische Methoden besprochen werden, wie die stetige Bewegung einer Maschine hervorgebracht werden kann.
Nachdem Leibniz im ersten Teil von N.\;\ref{57040} eine bereits 1674/1675 in Paris entworfene Windmühle, deren ununterbrochene Drehung durch die Erhebung hängender Gewichte gesichert werden soll, erneut dargelegt hat, widmet er sich den im Titel beider \textit{Schedae} angegebenen Messinstrumenten als möglichen mittelbaren Quellen einer stetigen Bewegung.
Der Kerngedanke hierbei ist, dass die täglichen Veränderungen dieser Messinstrumente, angemessen ,gespeichert‘, eine hinreichende Bewegungsquelle seien, um einen mechanischen Apparat -- etwa ein Uhrwerk -- stetig zu betätigen.
Der \textit{motus perpetuus}, von dem in diesen Texten die Rede ist, ist daher das Ergebnis einer Umwandlung mechanischer, von außen gespeister Kraft und stellt somit keine Verletzung des Erhaltungsgrundsatzes dar.
Aufgrund ihres einheitlichen Charakters werden die Aufzeichnungen N.\;\ref{57040} und \ref{57041} einheitlich ediert.
\pend%
\pstart%
Der Träger von N.\;\ref{57040} ist von Leibniz eigenhändig auf Oktober 1677 datiert, der Träger von N.\;\ref{57041} weist hingegen keine Datierung auf.
Da die Aufzeichnung N.\;\ref{57041} aber den Text von N.\;\ref{57040} unmittelbar fortsetzt, dürfte wohl auch die zweite Aufzeichnung unmittelbar oder kurz nach der ersten entstanden sein.
Demgemäß wird die Datierung von N.\;\ref{57040} auch für N.\;\ref{57041} übernommen.
Das gesamte Stück N.\;\ref{motus-perpetuus_intro} entstand mit großer Wahrscheinlichkeit im Oktober 1677.
\pend%
\selectlanguage{latin}%
\frenchspacing%
\normalsize%
